%%
% The BIThesis Template for Bachelor Graduation Thesis
%
% 北京理工大学毕业设计（论文）结论 —— 使用 XeLaTeX 编译
%
% Copyright 2020-2023 BITNP
%
% This work may be distributed and/or modified under the
% conditions of the LaTeX Project Public License, either version 1.3
% of this license or (at your option) any later version.
% The latest version of this license is in
%   http://www.latex-project.org/lppl.txt
% and version 1.3 or later is part of all distributions of LaTeX
% version 2005/12/01 or later.
%
% This work has the LPPL maintenance status `maintained'.
%
% The Current Maintainer of this work is Feng Kaiyu.
%
% Compile with: xelatex -> biber -> xelatex -> xelatex

\begin{conclusion}
语音情绪识别技术需要从非平稳语音信号中提取与情绪状态相关的声学线索。受用户差异、噪声干扰、语料分布变化和标注数据规模限制影响，模型在同源clean条件下取得的指标，未必能够直接代表实际应用环境中的稳定性。本文围绕对数梅尔频谱输入下的语音情绪识别方法开展研究，主要完成了数据集与信号处理流程整理、基础模型结构分析、卷积神经网络基线模型与Transformer系列模型对比，以及CNN-Conformer代表结构的进一步实验分析。

在数据与信号处理方面，本文选取RAVDESS语音情绪数据集作为主要实验对象，采用用户无关划分方式组织训练与评价。本文将原始语音统一重采样到16kHz，并采用幅值归一化、短时傅里叶变换、梅尔滤波器组映射、对数变换和特征标准化等处理，形成以对数梅尔频谱为中心的输入表示。针对发话长度不一致的问题，本文结合固定尺寸整理和分块组织方式，使不同模型能够在相对统一的声学输入基础上进行比较。

在模型结构方面，本文首先分析了前馈神经网络、卷积神经网络、循环神经网络、注意力机制和Transformer结构的基本特点。随后，本文实现并比较了三类模型：卷积神经网络基线模型、纯Transformer模型和CNN-Conformer模型。实验结果显示，卷积神经网络基线模型取得0.62196的Macro-F1、0.61667的Accuracy和0.61563的UAR，说明局部时频模式提取在当前任务条件下仍具有较强的参照意义。纯Transformer模型取得0.51163的Macro-F1、0.52000的Accuracy和0.51250的UAR，说明缺少显式局部归纳偏置时，纯注意力结构在当前数据规模下较难稳定转化为验证集收益。

CNN-Conformer模型在本文实验中取得较高的综合指标。代表配置采用无卷积前端时间分块输入、LayerNorm、192维嵌入、4层Conformer编码器、8个注意力头、768维前馈网络、卷积核大小31、末层输出、注意力池化和频谱级混合增强。该配置取得0.70563的Macro-F1、0.70000的Accuracy和0.70938的UAR。实验过程表明，CNN-Conformer模型的结果不主要来自结构规模的持续增加，该结果还与前端标记化方式、模型容量控制、池化策略和训练阶段正则化之间的组合有关。无卷积前端时间分块输入减少了多级卷积压缩对细粒度时间线索的影响，混合增强在一定程度上缓解了用户无关划分下的过拟合风险。

本文还围绕已选定的CNN-Conformer代表模型开展了噪声条件实验和跨语料库实验。噪声条件实验表明，模型在clean条件下的较高指标不能直接等同于受扰输入下的稳定表现；不同噪声类型和不同SNR区间会形成不同的退化轨迹，其中持续低频覆盖型噪声对当前结构影响较为明显。跨语料库实验以RAVDESS六类子集为源语料、CREMA-D六类子集为目标语料，在不使用目标语料监督信息的source-only设置下进行评价。源验证集Macro-F1为0.56897，而目标语料Macro-F1降至0.09243，说明当前模型对语料分布变化仍较为敏感。上述补充实验为本文结果提供了边界说明，也为后续鲁棒性建模提供了依据。

本文工作仍存在若干限制。第一，主要实验基于RAVDESS语音子集，数据规模和语料类型较为有限。第二，跨语料库实验只采用source-only基线设置，尚未引入域适配或域不变表示学习。第三，本文主要采用对数梅尔频谱作为输入表示，尚未系统比较自监督语音特征、多模态信息或文本内容对语音情绪识别的补充作用。后续研究可围绕真实噪声条件下的增强训练、跨语料库域适配、用户差异抑制、自监督语音表示轻量化应用以及多模态情绪识别展开，从而进一步提高模型在实际环境中的适应能力。
\end{conclusion}
